\FloatBarrier
\section{Exchange-Only Qubit Control}
\label{app:exchange-only-control}

Pair exchange acts along different logical axes within a three-spin
exchange-only qubit. A coupling between two encoded blocks additionally
allows amplitude to leave their individual code spaces, so an encoded
two-qubit gate must control the complete six-spin evolution. This
appendix gives the single-pair rotations, their calibration and the
conditions on a complete two-qubit sequence.

The logical basis is defined in Eq.~\eqref{eq:exchange-only-basis};
preparation and pair readout are developed in
Appendices~\ref{app:state-preparation} and~\ref{app:spin-readout}.

\subsection{Single-qubit gates for exchange-only qubits}
\label{sec:calibration-eo-rotations}
\label{sec:measurement-eo-exchange}
The two interdot barriers provide rotations about different logical
axes. We first identify those axes and the conditions under which a
single-pair pulse has a simple propagator. Pair readout then
calibrates the two pulse areas in different ways. Finally, we ask
whether simultaneous exchange can produce a direct logical $x$
rotation over the complete voltage waveform.

\subsubsection{Logical Hamiltonian and single-pair rotations}
\label{sec:calibration-eo-logical-axes}

In the three-spin basis of Eq.~\eqref{eq:exchange-only-basis},
logical 0 has pair 1--2 in a singlet and logical 1 has that pair
in a triplet. Exchange on 1--2 therefore produces a relative
phase in this basis. Exchange on 2--3 is diagonal in a different
pair-spin basis, so it both mixes and shifts the chosen logical
states. To establish the restriction to this two-dimensional
space, consider the physical Hamiltonian in
Eq.~\eqref{eq:exchange-only-control}. Each pair exchange is
$\mathbf{s}_i\cdot\mathbf{s}_j
=[(\mathbf{s}_i+\mathbf{s}_j)^2-3I/2]/2$,
so it commutes with all components of the pair's total spin and
with the remaining spin. It therefore commutes with
$\mathbf{s}_{\mathrm{tot}}^2$ and $s_{\mathrm{tot},z}$.
Since the code spans the full $(S,m)=(1/2,+1/2)$ sector,
arbitrary time-dependent exchange pulses stay within it.
Additional nonuniform magnetic fields need not preserve this space.

Define $P_{\mathrm E}=\ket{0_{\mathrm E}}\bra{0_{\mathrm E}}
+\ket{1_{\mathrm E}}\bra{1_{\mathrm E}}$.
Then $(I-P_{\mathrm E})H_{\mathrm E}P_{\mathrm E}=0$.
Restricting $P_{\mathrm E}H_{\mathrm E}P_{\mathrm E}$ means computing
its matrix elements in the encoded basis. Using
$\mathbf{s}_i\cdot\mathbf{s}_j=\mathrm{SWAP}_{ij}/2-I/4$ gives
\begin{equation}
 \begin{aligned}
 \left.P_{\mathrm E}\mathbf{s}_1\cdot\mathbf{s}_2P_{\mathrm E}
       \right|_{\mathrm E}
   &=\begin{pmatrix}-3/4&0\\0&1/4\end{pmatrix}
     =-\frac14 I-\frac12 Z_{\mathrm E},\\
 \left.P_{\mathrm E}\mathbf{s}_2\cdot\mathbf{s}_3P_{\mathrm E}
       \right|_{\mathrm E}
   &=\begin{pmatrix}0&\sqrt3/4\\\sqrt3/4&-1/2\end{pmatrix}
     =-\frac14 I+\frac14 Z_{\mathrm E}
                       +\frac{\sqrt3}{4}X_{\mathrm E}.
 \end{aligned}
 \label{eq:exchange-only-projection}
\end{equation}
These are the matrices of the physical exchanges acting within the
code, with $X_{\mathrm E},Z_{\mathrm E}$ the encoded Pauli operators.
The two compensated interdot barriers control $J_{12}$ and $J_{23}$.
Combining the matrices with these coefficients and dropping the
common energy $-(J_{12}+J_{23})I/4$ gives
\begin{equation}
 H_{\mathrm{E,q}}(t)
 =\frac{\sqrt3J_{23}(t)}4 X_{\mathrm E}
  +\left[-\frac{J_{12}(t)}2+\frac{J_{23}(t)}4\right]Z_{\mathrm E}.
 \label{eq:calibration-eo-axis}
\end{equation}
The following construction assumes selective exchange pulses.
To realize $U_{12,\mathrm E}$, pulse the 1--2 barrier so that
$J_{12}(t)$ is active while $J_{23}(t)$ is negligible. To realize
$U_{23,\mathrm E}$, pulse the 2--3 barrier with $J_{12}(t)$
negligible. Having two barriers does not by itself establish that
the inactive exchange can be suppressed; this is a condition to
check at the operating point. For either single-pair pulse, define
its exchange area by
\[
 \theta_{ij}=\frac1\hbar\int J_{ij}(t)\,dt.
\]
To derive the ideal propagators, set the inactive coupling to
zero in Eq.~\eqref{eq:calibration-eo-axis}. A 1--2 pulse then has
$H_{\mathrm{E,q}}(t)=-J_{12}(t)Z_{\mathrm E}/2$, whereas a 2--3
pulse has
$H_{\mathrm{E,q}}(t)=J_{23}(t)(\sqrt3X_{\mathrm E}+Z_{\mathrm E})/4$.
Integrating each of these single-pair Hamiltonians gives, up to
the global phase of the omitted identity term,
\begin{equation}
 \begin{aligned}
 U_{12,\mathrm E}(\theta_{12})
  &=\exp(i\theta_{12}Z_{\mathrm E}/2),\\
 U_{23,\mathrm E}(\theta_{23})
  &=\exp\!\left[-\frac{i\theta_{23}}4
       (\sqrt3X_{\mathrm E}+Z_{\mathrm E})\right].
 \end{aligned}
 \label{eq:calibration-eo-propagators}
\end{equation}
If the inactive coupling is appreciable, the actual pulse is not
described by these single-pair propagators; its evolution must be
calculated from Eq.~\eqref{eq:calibration-eo-axis} with both couplings.
These are the two nonparallel logical rotations summarized earlier
in Eq.~\eqref{eq:exchange-only-rotations}; they are written here
because the following measurements calibrate their two pulse areas.

\subsubsection{Calibrating the two pair-exchange pulse areas}
\label{sec:calibration-eo-pulse-areas}

Pair readout distinguishes singlet and triplet character on dots
1--2. A 2--3 pulse mixes those logical states, so its area appears
directly as a population oscillation. A 1--2 pulse changes only
their relative phase and needs a preparation and analysis rotation
before the same readout can reveal it.

\paragraph{Measuring the 2--3 coupling.}
Prepare $\ket{0_{\mathrm E}}$ using
Section~\ref{sec:calibration-eo-initialization}, suppress $J_{12}$,
and pulse only the 2--3 barrier. For a constant $J_{23}=J$,
\begin{equation}
 H_{23}^{\mathrm E}=\frac J4Z_{\mathrm E}
                   +\frac{\sqrt3J}{4}X_{\mathrm E},\qquad
 P_{1_{\mathrm E}}(\tau)=\frac34
                    \sin^2\!\left(\frac{J\tau}{2\hbar}\right).
 \label{eq:measurement-eo-oscillations}
\end{equation}
The oscillation reaches $3/4$ rather than 1 because this rotation
axis is tilted relative to the preparation and pair-readout axis.
Vary the hold time, return to the pair-readout setting on dots
1--2, and count pair-triplet outcomes using
Section~\ref{sec:calibration-eo-readout}. The period gives $J/h$.
Repeat at different compensated barrier voltages to determine
$J_{23}$ along the control path. Pair-triplet leakage states can
produce the same blocked record, so this interpretation assumes
that the evolution remains in the code space.

For a pulse with variable coupling but $J_{12}=0$, the same
probability is $3\sin^2(\theta_{23}/2)/4$, where
$\theta_{23}=\hbar^{-1}\int J_{23}(t)\,dt$.
A duration scan with fixed edges calibrates this complete pulse
area, including its edge contribution.

\paragraph{Measuring the 1--2 phase.}
A 1--2 pulse only adds phase to the initial
$\ket{0_{\mathrm E}}$, so first apply
$A=U_{23,\mathrm E}(\beta)$ with $\cos\beta=-1/3$.
The expression above gives
$P_{1_{\mathrm E}}=3(1-\cos\beta)/8=1/2$:
$A\ket{0_{\mathrm E}}$ is an equal-weight superposition.
Insert the 1--2 pulse and then apply $A^\dagger$ before pair readout.
Using the pair propagators in
Eq.~\eqref{eq:calibration-eo-propagators},
\begin{equation}
 P_{0_{\mathrm E}}=
 \left|\bra{0_{\mathrm E}}A^\dagger
       e^{i\theta_{12}Z_{\mathrm E}/2}A
       \ket{0_{\mathrm E}}\right|^2
 =\frac{1+\cos\theta_{12}}2.
 \label{eq:calibration-eo-phase-sequence}
\end{equation}
The inverse rotation can be supplied by a positive 2--3 pulse of
area $2\pi-\beta$, up to global phase. Vary the 1--2 hold time
between these fixed preparation and analysis pulses to calibrate
$\theta_{12}$ from the return oscillation.
Residual couplings during either pulse must be included in the
two-axis Hamiltonian above.

\subsubsection{Testing a direct logical \texorpdfstring{$x$}{x} rotation}
\label{sec:calibration-eo-direct-x}

First consider constant simultaneous couplings $J_{12}$ and
$J_{23}$ during a hold of duration $\tau$. With $J_{23}=0$,
the hold gives a pure $z$ rotation $R_z(-\theta_{12})$.
With $J_{12}=J_{23}/2$, the $Z_{\mathrm E}$ coefficient
cancels and the hold gives a pure $x$ rotation:
\begin{equation}
 H_{\mathrm{E,q}}=\frac{\sqrt3J_{23}}4X_{\mathrm E},\qquad
 \vartheta_x=\frac{\sqrt3J_{23}\tau}{2\hbar}.
 \label{eq:calibration-eo-pure-x}
\end{equation}
For arbitrary constant values of the two couplings, pair readout
after the hold gives
\begin{equation}
 P_{1_{\mathrm E}}(\tau)=
 \frac{3J_{23}^2}{3J_{23}^2+(J_{23}-2J_{12})^2}
 \sin^2\left(
  \frac{\tau\sqrt{3J_{23}^2+(J_{23}-2J_{12})^2}}{4\hbar}
       \right).
 \label{eq:calibration-eo-axis-tuning}
\end{equation}
At fixed $J_{23}$, scanning $J_{12}$ and $\tau$ identifies the
constant setting with vanishing $z$ component by its maximum
oscillation contrast; the period then determines the $x$-rotation
rate. The two barrier voltages must be calibrated when applied
together. Their separately measured responses only suggest an
initial setting, and readout contrast and decay must be accounted
for when interpreting the measured amplitude.

A real voltage pulse can have time-dependent couplings on its
edges. The constant-hold result does not by itself show that the
complete pulse is a pure $R_x$. One sufficient condition is
$J_{12}(t)=J_{23}(t)/2$ throughout the active waveform, with
other logical terms negligible; then
$\vartheta_x=(\sqrt3/2\hbar)\int J_{23}(t)\,dt$.
If the ratio changes, the complete evolution must be evaluated
and calibrated as a gate. When the available waveform does not
produce the desired $R_x$, compose it from the two calibrated
nonparallel pair rotations instead.

These pulse-area and axis measurements characterize coherent gates.
An encoded-qubit relaxation measurement asks a separate question:
prepare an excited logical state and measure its population after
a wait, as in Section~\ref{sec:measurement-relaxation}. Extra
leakage channels may require tracking more than two populations.

\subsection{Two-qubit gates for exchange-only qubits}
\label{sec:calibration-eo-two-qubit-gates}

Let qubit $A$ occupy spins $(1,2,3)$ and qubit $B$ occupy
$(4,5,6)$. Exchange within each triple dot produces the one-qubit
rotations of the preceding section. A two-qubit gate additionally
uses the boundary coupling $J_{34}$ between the two blocks. That
coupling can entangle them, but it does not preserve each block's
code space by itself. The control problem is therefore to combine
intra-block and inter-block exchange pulses so that the final state
returns to both code spaces while retaining a nonlocal logical
operation.

\subsubsection{Why one boundary-exchange pulse is insufficient}

For a linear six-dot array, the physical Hamiltonian is
\begin{equation}
 H_6(t)=\sum_{i=1}^{5}J_{i,i+1}(t)\,
                     \mathbf s_i\cdot\mathbf s_{i+1}.
 \label{eq:calibration-six-spin-hamiltonian}
\end{equation}
The boundary term $J_{34}\mathbf s_3\cdot\mathbf s_4$ transfers
spin angular momentum between the blocks. If
$\mathbf S_A=\mathbf s_1+\mathbf s_2+\mathbf s_3$ and
$\mathbf S_B=\mathbf s_4+\mathbf s_5+\mathbf s_6$, then
\[
 [\mathbf s_3\cdot\mathbf s_4,\mathbf S_A^2]\ne0,
 \qquad
 [\mathbf s_3\cdot\mathbf s_4,\mathbf S_B^2]\ne0.
\]
A single boundary pulse can therefore move amplitude from an
$S=1/2$ encoded block into an $S=3/2$ state. It may entangle the
six physical spins without defining a unitary gate on the two
logical qubits.

\subsubsection{Conditions on a complete exchange sequence}

A candidate encoded gate is an ordered product
\begin{equation}
 U_6=U_{i_Mj_M}(\theta_M)\cdots U_{i_1j_1}(\theta_1),
 \qquad
 U_{ij}(\theta)=e^{-i\theta\,\mathbf s_i\cdot\mathbf s_j},
 \qquad
 \theta=\frac1\hbar\int J_{ij}(t)\,dt .
 \label{eq:calibration-six-spin-sequence}
\end{equation}
The order matters because exchanges on overlapping pairs do not
commute. Let $P_{AB}$ project onto the product of the two
exchange-only code spaces. A valid encoded gate must satisfy
\begin{equation}
 (I-P_{AB})U_6P_{AB}=0,
 \qquad
 \left.P_{AB}U_6P_{AB}\right|_{\mathrm{code}}
   =e^{i\chi}U_{\mathrm{logical}}.
 \label{eq:calibration-encoded-gate-conditions}
\end{equation}
The first condition allows intermediate leakage but requires its
amplitude to cancel at the end. The second condition fixes the
logical action after that return. For an encoded CNOT,
$U_{\mathrm{logical}}=\mathrm{CNOT}_{A\to B}$, reproducing the
target stated in Eq.~\eqref{eq:encoded-cnot-target}.

\subsubsection{Calibrating and testing the complete sequence}

Each pair pulse is first calibrated through its measured exchange
area, but those separate calibrations do not establish the
six-spin gate. Program the ordered pulse list, prepare each of the
four logical basis inputs, and apply the complete sequence before
the exchange-only readout of both blocks. The final code-space
population tests the first condition above; the logical outcomes
test the basis-state action of the second. Superposition inputs and
analysis rotations are additionally required to determine the
relative phases and distinguish the target gate from a truth table
with incorrect coherent phases.

Exact leakage-free exchange-only CNOT sequences supply a concrete
pulse ordering and set of exchange areas~\cite{FongWandzura2011}.
Their individual areas are initial settings, not experimental
facts: the full sequence must be retested after the pair responses,
switching edges and residual simultaneous couplings are included.
A single calibrated $J_{34}$ pulse cannot replace this complete
sequence.
